% !PRINTED none
% Chapter-opening page: centred "Zweiter Abschnitt." / "Theta-Funktionen." with a rule
% beneath, then the \S\,15 heading and running text. No page number is printed on
% this page (chapter-opening convention).

\begin{center}
\gesperrt{Second Section.}

\medskip
\textbf{\gesperrt{Theta Functions.}}
\end{center}

\bigskip
\begin{center}\rule{2cm}{0.4pt}\end{center}
\bigskip

\begin{center}
\textbf{\S\,15. \gesperrt{Assumptions from the Theory of Functions.}}
\end{center}

The concept of doubly periodic functions can be rightly grasped only if one regards them as functions of a complex argument $u = v + iw$, in which $i$ has, as usual, the meaning of $\sqrt{-1}$. The variable $u$ counts for us as an independent and unrestrictedly variable quantity, which, following the precedent of \gesperrt{Gauss}, is geometrically represented by the points of a plane in such a way that the point whose coordinates in a rectangular system are $v, w$ is regarded as the carrier of the value $u = v + iw$ and is designated briefly as the point $u$.

We here restrict our consideration to single-valued analytic functions $\varphi(u)$ which have no \gesperrt{essentially singular} points in the finite region. What happens at infinity we leave undecided. We therefore assume the following of the functions under consideration:

1. In every finite portion of surface there lies a finite number of points at which the function $\varphi(u)$ becomes infinite; these points are called \gesperrt{points of discontinuity}.

The number of points of discontinuity altogether, i.\,e. in the whole infinite plane, may of course also be infinitely great.

2. If $u_0$ is a point not belonging to the points of discontinuity, then the function is developable in a series proceeding by integral ascending positive powers of $u - u_0$, which is convergent in a circle having the point $u_0$ as its centre and reaching as far as the nearest point of discontinuity. One says that the function has, in the neighbourhood of the point $u_0$, the \gesperrt{character of an integral function}.
